\cotefile{1}
\scanpage{4}
Soit $C$ une $C^*$-algèbre, on veut déterminer les \deleted{\unsure{endomorphismes}} endomorphismes linéaires continus $T$ permutant aux translations gauches. Soit $T$ un \unsure{telle|tel}, son bitransposé $T''$ est un endomorphisme \unsure{fblt} continu \inserted{de $C''$} qui permute aux translations $x\unsure{\mapsto|\to}xy$ \deleted{$(y\in C)$ donc} pour $y\in C$, donc pour \deleted{tout} $y\in C''$ (\inserted{donc} $T''(xy)=T''(x)y$ \deleted{\unsure{il est} \illegible{2} \unsure{donc}} \inserted{les deux} mbs sont \unsure{fblt} continus). Donc de la forme $x\to ux$ ($u\in C''$ \inserted{bien déterminé de \unsure{norme} $\unsure{\lVert T''\rVert|\lVert T\rVert}$}). En fait, $u$ appartient à l'ens. $\Gamma_g$ des éléments $u$ de $C''$ tels que $uC\subset C$. On voit donc que l'ensemble d'endomorphismes cherché est isomorphe, avec sa structure d'algèbre normée, à $\Gamma_g$, \deleted{\unsure{l'hermitien}} $u$ hermitien signifie $x^*T(x)$ hermitien pour \unsure{tout} $x\in C$ ; $u$ positif signifie $x^*T(x)\geq 0$ pour \deleted{tout} $\unsure{u|x}\in C$.

Soit \deleted{$\unsure{\mathfrak g|g}$} $\Gamma$ l'ens des $u\in C''$ tels que $u\underline{C}\subset\underline{C}$, $\underline{C}u\subset\underline{C}$. C'est une sous $C^*$-algèbre \unsure{involutive}. \unsure{Tout} \unsure{est|elt} \unsure{alors|élément} hermitien de $\Gamma_g$ est dans $\Gamma$

Supposons maintenant $C$ \unsure{plongé|prolongé} dans la $C^*$-algèbre $R$, et soit $\Gamma_1$ l'ens des $x\in R$ tels que $xC\subset C$, $Cx\subset \unsure{R|C}$, c'est une sous $C^*$-algèbre de $R$. Pour tout $a\in\Gamma_1$, soit $\varphi(a)$ l'élément de $C''$ correspondant à l'opération $T_a(x)=ax$ dans $C$. Si $a$ hermitien, $\varphi(a)$ est hermitien, donc dans $\Gamma$, donc pour \deleted{tout} $a\in\Gamma_{\unsure{g|1}}$, $\varphi(a)\in\Gamma$. De plus, $\varphi$ est évidemment une représentation d'algèbre, compatible avec l'involution (puisque
% archive sheet number: 5

